\chapter{Topología del grafo: treewidth y curvatura}
\label{cap:topologia}

Hasta acá los invariantes del TypeGraph fueron \emph{semánticos}: link
counts, monotonía de capabilities, ausencia de ciclos de directorios.
Este capítulo introduce dos métricas \emph{estructurales} que el
runtime computa sobre el grafo: \emph{treewidth} (aproximada) y
\emph{curvatura de Ollivier-Ricci}. Ambas miden qué tan ``árbol-like''
es el grafo. La motivación es operativa: si un workload adversarial
empuja el grafo lejos del régimen tree-like, se levantan alertas; si
una proyección del grafo (cap.~\ref{cap:decepcion}) tiene curvatura
sistemáticamente distinta del original, el atacante podría detectarla
por timing-side-channel.

\section{Por qué importan estas métricas en un FS}

\paragraph{Treewidth.} Una cantidad asociada a cómo de bien un grafo se
puede descomponer en un árbol de bags pequeños. Treewidth bajo
$\Leftrightarrow$ grafo tree-like $\Leftrightarrow$ muchos algoritmos
NP-difíciles se vuelven polinomiales sobre él. Para sotFS, un
treewidth acotado superiormente por el \texttt{LINK\_MAX = 65535}
(cap.~\ref{cap:typegraph}) más unas pocas unidades garantiza
escalabilidad de algoritmos de análisis y monitoreo.

\paragraph{Curvatura.} La curvatura de Ollivier-Ricci sobre un grafo
mide qué tan ``positivamente curvado'' (cliques, esferas) o
``negativamente curvado'' (árboles, cuellos de botella) es cada arista.
Permite detectar patrones adversariales: una ``symlink bomb''
produce una firma de curvatura distinta a la de un FS normal; una
operación de \texttt{mkdir} masivo produce una spike positiva en el
directorio padre.

\section{Treewidth: definición y aproximación}

\begin{defn}[Tree decomposition]
\label{def:tree-decomposition}
Una \emph{tree decomposition} de un grafo $G = (V, E)$ es un árbol $T$
junto con una asignación $X: V_T \to \mathcal{P}(V)$ tal que:
\begin{enumerate}
  \item $\bigcup_{t \in V_T} X(t) = V$.
  \item Para toda arista $(u, v) \in E$, existe $t$ con
    $\{u, v\} \subseteq X(t)$.
  \item Para todo $v \in V$, el conjunto $\{t : v \in X(t)\}$ induce
    un subárbol conexo de $T$.
\end{enumerate}
El \emph{ancho} de la descomposición es $\max_t |X(t)| - 1$. El
\emph{treewidth} de $G$ es el mínimo ancho sobre todas las
descomposiciones.
\end{defn}

Calcular treewidth exacto es NP-hard \cite{Bodlaender1993}. sotFS usa
una aproximación greedy con \emph{min-degree elimination ordering}.

\subsection*{Algoritmo greedy}

\begin{algorithm}[H]
\SetAlgoLined
\KwIn{Grafo $G = (V, E)$}
\KwOut{Upper bound de treewidth}
\caption{\textsc{GreedyMinDegreeTreewidth}}\label{alg:treewidth}

$G' \gets G$
$\mathit{tw} \gets 0$\;
\While{$|V_{G'}| > 1$}{
  $v \gets \arg\min_{u \in V_{G'}} \deg_{G'}(u)$\;
  $\mathit{tw} \gets \max(\mathit{tw}, \deg_{G'}(v))$\;
  
  \ForEach{$u_1, u_2 \in N_{G'}(v), u_1 \neq u_2$}{
    Add edge $(u_1, u_2)$ to $G'$ if not present\;
  }
  Remove $v$ from $G'$\;
}
\Return{$\mathit{tw}$}\;
\end{algorithm}

\paragraph{Costo.} $O(n^2)$ con representación adecuada. Suficientemente
barato para correr periódicamente sobre el grafo del FS.

\paragraph{Resultado para sotFS.}

\begin{teorema}[Treewidth acotado para FS con hard-link cap]
\label{teo:treewidth-fs}
Sea $G$ un TypeGraph well-formed con $i.\texttt{link\_count} \leq
\texttt{LINK\_MAX}$ para todo inode. Entonces el treewidth de $G$
satisface $\mathrm{tw}(G) \leq \texttt{LINK\_MAX} + O(1)$.
\end{teorema}

Sin hard links ($\texttt{link\_count} \leq 1$ para todo inode), el
TypeGraph es esencialmente un árbol, $\mathrm{tw}(G) \leq 6$. La cota
\texttt{LINK\_MAX} de \invref{inv:link-count} (vía
\trampref{tr:link-treewidth}) preserva este orden de magnitud
incluso bajo workload adversarial.

\section{Curvatura de Ollivier-Ricci}

\begin{defn}[Curvatura discreta]
\label{def:ricci}
La \emph{curvatura de Ollivier-Ricci} de una arista $(u, v)$ en un
grafo $G$ es
\[
  \kappa(u, v) \;=\; 1 - \frac{W_1(\mu_u, \mu_v)}{d(u, v)}
\]
donde $W_1$ es la distancia Wasserstein-1, $d(u,v)$ es la distancia
en el grafo, y $\mu_u, \mu_v$ son medidas de probabilidad sobre los
vecinos de $u$ y $v$ respectivamente.
\end{defn}

Para grafos no ponderados, $d(u, v) = 1$ (los extremos son adyacentes),
y se usa la \emph{lazy random walk} con parámetro $\alpha \in [0, 1]$:
$\mu_x(x) = \alpha$, $\mu_x(z) = (1-\alpha)/\deg(x)$ para todo $z$
vecino de $x$.

\subsection*{Aproximación Lin-Lu-Yau}

Calcular $W_1$ exactamente requiere resolver un transporte óptimo
(linear program). \citaCrate{sotfs-monitor} usa la aproximación
Lin-Lu-Yau \cite{LinLuYau2011} con $\alpha = 0.5$:

\[
  \kappa_{\mathrm{LLY}}(u, v) \;\approx\;
    \frac{|N(u) \cap N(v)|}{\max(\deg u, \deg v)}
    \cdot (1-\alpha)^2 + \text{(correcciones de borde)}
\]

\paragraph{Costo.} $O(\deg u + \deg v)$ por arista, $O(E \cdot \bar{d})$
total. Implementación en \citaLine{sotfs-monitor/src/curvature.rs}{75}.

\subsection*{Algoritmo}

\begin{algorithm}[H]
\SetAlgoLined
\KwIn{Grafo $G$, arista $(u, v)$, parámetro $\alpha$}
\KwOut{Aproximación de $\kappa_{\mathrm{LLY}}(u, v)$}
\caption{\textsc{OllivierRicciLLY}}\label{alg:curvatura}

$\mathit{common} \gets |N(u) \cap N(v)|$\;
$\mathit{d}_u \gets \deg u; \;\; \mathit{d}_v \gets \deg v$\;
$\mathit{base} \gets \mathit{common} / \max(\mathit{d}_u, \mathit{d}_v)$\;
$\mathit{self\_transport} \gets (1 - \alpha)^2$\;
$\kappa \gets 1 - (1 - \mathit{base} \cdot \mathit{self\_transport})$\;
\Return{$\kappa$}\;
\end{algorithm}

\paragraph{Interpretación.}

\begin{itemize}
  \item $\kappa = 1$: arista en una clique (los vecindarios coinciden).
  \item $\kappa = 0$: arista en un grafo regular ``neutro''.
  \item $\kappa < 0$: arista en una región tree-like (cuellos de
    botella, dirs profundos).
\end{itemize}

\section{Detección de anomalías}

\begin{defn}[Anomalía por curvatura]
\label{def:anomalia}
Una arista $(u, v)$ es \emph{anómala} si
\[
  |\kappa(u, v) - \bar{\kappa}| > 2\,\sigma_{\kappa}
\]
donde $\bar{\kappa}$ y $\sigma_{\kappa}$ son la media y desviación
estándar de la curvatura sobre todas las aristas del grafo.
\end{defn}

\paragraph{Firmas observadas} (en
\citaLine{sotfs-monitor/src/curvature.rs}{173}):

\begin{itemize}
  \item \emph{Mass file creation}: spike positivo de $\kappa$ en el
    directorio padre (sus vecindarios crecen mucho).
  \item \emph{Symlink bomb}: $\kappa$ fuerte positivo en aristas
    \texttt{Symlink → target}.
  \item \emph{Directorio profundo}: $\kappa \approx -1.0$ en el cuello
    de la cadena.
\end{itemize}

\section{El spec TLA de curvatura}

\citaTLA{sotfs\_curvature.tla} (301 LOC) modela la curvatura como
una función computada del estado y prueba dos cosas:

\begin{lstlisting}[language=tla, caption={Invariantes de curvatura
en TLA.}, label={lst:curvature-tla}]
SafetyInvariant ==
    /\ TypeInvariant
    /\ CurvatureBounded
    /\ KappaConsistent
    /\ CurvatureDropImpliesFanOut
\end{lstlisting}

\begin{description}
  \item[\texttt{CurvatureBounded}] $|\kappa(u,v)| \leq \max(\deg u,
    \deg v)$.
  \item[\texttt{KappaConsistent}] el cálculo es determinista: dos
    invocaciones sobre el mismo estado dan el mismo resultado.
  \item[\texttt{CurvatureDropImpliesFanOut}] si entre dos estados
    $\Delta\kappa < -T$ para alguna arista, entonces \emph{en el mismo
    paso} ocurrió un evento de fan-out (un nodo aumentó su grado en
    $\geq$ \texttt{FanOutDelta}).
\end{description}

TLC verifica que estas tres se mantienen bajo todas las acciones del
spec, incluyendo creates masivos y unlinks.

\section*{Resumen del capítulo}

\begin{itemize}
  \item \emph{Treewidth} mide qué tan tree-like es el grafo; sotFS lo
    aproxima con min-degree elimination ordering en $O(n^2)$
    (\citaLine{sotfs-monitor/src/treewidth.rs}{33}).
  \item Con \texttt{LINK\_MAX = 65535}, $\mathrm{tw}(G) \leq
    \texttt{LINK\_MAX} + O(1)$; sin hard links, $\leq 6$.
  \item \emph{Curvatura de Ollivier-Ricci} mide cuán ``positivamente
    curvada'' es cada arista (Wasserstein-1 entre vecindarios). sotFS
    usa la aproximación Lin-Lu-Yau con $\alpha = 0.5$, costo
    $O(\deg u + \deg v)$ por arista.
  \item Anomalías: aristas con $|\kappa - \bar{\kappa}| > 2\sigma$.
    Firmas distintivas para mass-create, symlink bomb, dir profundo.
  \item \citaTLA{sotfs\_curvature.tla} verifica tres invariantes:
    bounded, consistent, drop-implies-fan-out.
\end{itemize}

\section*{Ejercicios}

\begin{ejercicio}[\dificultad{$\star$}]
\label{ej:topo-1}
Defina treewidth en una línea. ¿Para qué algoritmos sirve tenerlo
acotado?
\end{ejercicio}

\begin{ejercicio}[\dificultad{$\star$}]
\label{ej:topo-2}
¿Cuál es el valor de $\alpha$ usado por sotFS en la lazy random walk
de Ollivier-Ricci?
\end{ejercicio}

\begin{ejercicio}[\dificultad{$\star\star$}]
\label{ej:topo-3}
Calcule la curvatura aproximada (Lin-Lu-Yau, $\alpha = 0.5$) de una
arista en un grafo triangular (vecindarios coinciden completamente).
\end{ejercicio}

\begin{ejercicio}[\dificultad{$\star\star$}]
\label{ej:topo-4}
Calcule la curvatura aproximada de una arista que une dos hojas de un
árbol binario completo. ¿Por qué da negativa?
\end{ejercicio}

\begin{ejercicio}[\dificultad{$\star\star$}]
\label{ej:topo-5}
Demuestre el \teoref{teo:treewidth-fs} para el caso sin hard links
(\texttt{link\_count} $\leq$ 1 para todos los inodes). Pista: identificá
los seis tipos de nodo como candidatos a bag.
\end{ejercicio}

\begin{ejercicio}[\dificultad{$\star\star$}]
\label{ej:topo-6}
Considere un workload adversarial que crea $10^4$ archivos en un solo
directorio. ¿Qué firma de curvatura esperamos en el inode del
directorio? ¿En las aristas \econtains{} salientes?
\end{ejercicio}

\begin{ejercicio}[\dificultad{$\star\star\star$}]
\label{ej:topo-7}
\citaTLA{sotfs\_curvature.tla} prueba que un drop de curvatura
implica un fan-out en el mismo paso. ¿Vale la implicación inversa
(fan-out implica drop)? Justifique.
\end{ejercicio}

\begin{ejercicio}[\dificultad{$\star\star\star$}]
\label{ej:topo-8}
Diseñe un detector de anomalías que use \emph{ambos} treewidth y
curvatura. ¿Cómo combinás las dos señales? Pista: un cambio brusco de
treewidth con $\kappa$ estable indica algo distinto a un $\kappa$
shift con treewidth estable.
\end{ejercicio}
